\begin{proof}
We first compare the two fieldwise constructions in
\cref{lem:p171-module-c-field-comparison}.  We use the same ordered
intersection line on both sides.  \cite[Section~8.3 and Appendix~A]{BE21}
defines the Deligne pairing by the determinant finite difference used in
\cite{Ducrot97CubeIntersection}; \cite[Theorem~A.22]{BE21} gives its
positive restriction isomorphism, and \cite[Example~8.15]{BE21} identifies
its dimension-zero member with the positive finite norm.  Thus the
evaluation element below is not silently replaced by an element of the dual
line.

We need the following comparison, which we prove rather than cite as a
change of terminology.  The recursive calculation is first made after the
complete isometric scalar extension of
\cite[Definition~3.6.1]{ChenMoriwaki21AdelicIntersection}, so that
\cite[Remark~3.7.3 and formula~(3.16)]{ChenMoriwaki21AdelicIntersection}
applies.  Descent to the original nontrivially valued field is justified
below by
\cite[Proposition~3.8.9 and Remark~3.8.10]{ChenMoriwaki21AdelicIntersection}
and by the corresponding Deligne base-change argument.  Suppose that
\(t=(t_0,\ldots,t_n)\), with \(t_i\in E_i\), is such that the first
\(n\) sections form a regular sequence on all of \(Z\) and the total
intersection is empty.  Let
\[
 e_t\in I_{Z/K}(L_0,\ldots,L_n)\setminus\{0\}
\]
be obtained by successive positive restriction and the final finite norm.
Then
\begin{equation}
 h^{t}_{\mathrm{CM}}(Z;\vartheta_0,\ldots,\vartheta_n)
 =\lambda_K(e_t;\vartheta_0,\ldots,\vartheta_n).
 \label{eq:p171-module-c-height-weight}
\end{equation}
For \(n=0\), over the complete algebraic closure,
\cite[formula~(3.15)]{ChenMoriwaki21AdelicIntersection} is
\(-\int\log\lVert t_0\rVert_{\vartheta_0}\,\delta_Z\).
\cite[Theorem~8.16, formula~(8.11)]{BE21} says that the logarithm of the norm
of \(e_t\) is the same integral with the positive sign.  Taking its negative
proves \eqref{eq:p171-module-c-height-weight} there.  The base-change
argument below descends the equality to \(K\), and
\cite[Example~8.2]{BE21} identifies the descended measure as
\[
 \delta_Z=\sum_z\operatorname{length}(A_z)[k_z:K]\,\delta_z
\]
for a finite scheme \(Z=\operatorname{Spec}\prod_zA_z\).  Thus the descended
base case retains every residue degree and every nonreduced Artin length.

Assume the equality in dimensions below \(n\).  The local intersection
number is symmetric by
\cite[Remark~3.6.2(1)]{ChenMoriwaki21AdelicIntersection}.  We may therefore
apply
\cite[Definition~3.7.1 and formula~(3.16)]{ChenMoriwaki21AdelicIntersection}
successively in the order
\(t_0,\ldots,t_{n-1}\).  At the \(j\)-th step, \(t_j\) is regular on the
entire current quotient, so its effective Cartier cycle is the top
fundamental cycle of that quotient, with the coefficients supplied by
Cartier intersection.  This recursion restricts the remaining
local height to that cycle and subtracts the integral of
\(\log\lVert t_j\rVert\) against the mixed measure in the other slots.
\cite[Proposition~8.5 and Theorem~8.16, formula~(8.8)]{BE21} applies the
same positive restriction to the Deligne evaluation element and subtracts
the same integral.  After \(n\) restrictions, \(t_n\) is nonvanishing on the
finite scheme and the already proved dimension-zero comparison applies.
Induction proves \eqref{eq:p171-module-c-height-weight}.  This also shows
directly that the comparison is linear in
\([Z]_n=\sum_T\ell(\mathcal O_{Z,\eta_T})[T]\).  The analytic measures in
\cite[formula~(8.1) and Proposition~8.5]{BE21} use that same cycle.
Thus a generic nilpotent thickness is counted by its length, while a
component of dimension below \(n\), whether embedded or isolated,
contributes no extra summand.

We also record the field-extension step, since normalization after an
unsupported extension would leave a gap.  If \(F/K\) is a complete
isometric valued extension, the algebraic Deligne line commutes with base
change.  The pullback of its metric satisfies the defining restriction
formula of \cite[Theorem~8.16]{BE21}:
\cite[Proposition~8.3(iv)]{BE21} pushes the mixed measure on \(Z_F\) to the
mixed measure on \(Z\), and algebraic restriction commutes with base change.
\cite[Theorem~8.16, formula~(8.11)]{BE21} gives the same assertion in
dimension zero.  Uniqueness in \cite[Theorem~8.16]{BE21} therefore gives
\begin{equation}
 \lambda_F(e\otimes1;\vartheta_{0,F},\ldots,\vartheta_{n,F})
 =\lambda_K(e;\vartheta_0,\ldots,\vartheta_n).
 \label{eq:p171-module-c-deligne-base-change}
\end{equation}
On the Chen--Moriwaki side,
\cite[Proposition~3.8.9]{ChenMoriwaki21AdelicIntersection} proves invariance
for non-Archimedean complete valued extensions, and
\cite[Remark~3.8.10]{ChenMoriwaki21AdelicIntersection} gives the Archimedean
case.  \cite[Definition~3.8.2 and Remark~3.8.3]{ChenMoriwaki21AdelicIntersection}
gives the corresponding scalar-extension identity for the epsilon polynomial
norm.  These are the required transport results if a normalizing tuple is
selected only after a complete isometric valued extension;
\cite[Remark~1.6.8]{ChenMoriwaki21AdelicIntersection} alone would not
suffice.
Consequently the recursively proved equality
\eqref{eq:p171-module-c-height-weight} descends from the complete algebraic
closure to every nontrivially valued complete field (and from \(\mathbf C\)
to \(\mathbf R\)).

Assume first that the valuation is nontrivial.  Then \(K\) is infinite.
Choose the sections successively.  Before choosing \(t_j\), form, for every
subset \(I\subset\{0,\ldots,j-1\}\), the scheme-theoretic intersection
\(Z_I=\bigcap_{i\in I}V(t_i)\), with \(Z_\varnothing=Z\), and avoid the
vanishing subspaces attached to every point in the finite union
\(\bigcup_I\operatorname{Ass}(Z_I)\).  Each forbidden set is a proper
linear subspace because the system is basepoint-free, and a finite union of
proper subspaces does not exhaust a vector space over an infinite field.
Thus every chosen section is a nonzerodivisor on every nonempty preceding
subset intersection.  In particular, each \(t_j\) is Cartier on \(Z\), the
family of divisors intersects properly, and
\(t_0,\ldots,t_{n-1}\) is a regular sequence on the full current
quotients.  After \(n\) cuts the full intersection has dimension at most
zero.  Choose \(t_n\) subject to the same associated-point conditions and
also nonvanishing at every point of that finite support.  The total
intersection is empty.  The same construction may be made after a complete
isometric valued extension; in that form descent uses
\cite[Proposition~3.8.9 and Remark~3.8.10]{ChenMoriwaki21AdelicIntersection}.

Put \(c_t=F_e(t)\).  Positive normalization of the line-valued resultant
is the pair of literal equalities
\begin{equation}
 e_t=e\,c_t,\qquad R_Z^{t}=F_e/c_t,
 \qquad R_Z^{t}(t)=1.
 \label{eq:p171-module-c-normalization-cancellation}
\end{equation}
The scalar \(c_t\) is nonzero.  For a nontrivially non-Archimedean
valuation, \cite[Theorem~3.8.7]{ChenMoriwaki21AdelicIntersection} and
\eqref{eq:p171-module-c-height-weight} give
\[
 \lambda_K(e_t)=\log\lVert R_Z^{t}\rVert_\varepsilon
 =\log\lVert F_e\rVert_\varepsilon-\log\lvert c_t\rvert.
\]
On the other hand, the frame law for the Deligne metric and the first
identity in \eqref{eq:p171-module-c-normalization-cancellation} give
\(\lambda_K(e_t)=\lambda_K(e)-\log\lvert c_t\rvert\).  The two scalar
terms cancel, proving \eqref{eq:p171-module-c-field-identity} with the
claimed sign.

For the usual complex absolute value,
\cite[Theorem~3.8.7]{ChenMoriwaki21AdelicIntersection} gives
\(h^t_{\mathrm{CM}}=\operatorname{ph}(R_Z^t)\), where
\(\operatorname{ph}\) is exactly the sphere integral plus the displayed
harmonic correction.  Its logarithmic homogeneity and
\eqref{eq:p171-module-c-normalization-cancellation} give the same
cancellation.  The real assertion follows from the scalar extension to
\(\mathbf C\) in \cite[Remark~3.8.10]{ChenMoriwaki21AdelicIntersection},
\cite[Remark~8.1]{BE21}, and
\eqref{eq:p171-module-c-deligne-base-change}; conjugation changes neither
the integral nor the resultant.  In particular, conjugate points are
pushed forward and not averaged, so no factor two occurs.

For the powered Archimedean field \(K_\rho\), unpower every norm and metric.
By \cref{thm:powered-normalization}, \eqref{eq:PN1-r3} identifies the mixed
measure with the usual unpowered measure, while \eqref{eq:PN2-r3}
identifies the additive Deligne weight with \(\rho\) times the usual
weight.  The dimension-zero anchor is \cref{lem:zero-anchor-r3}.
Multiplying the whole usual formula by \(\rho\), including the harmonic
correction, proves \eqref{eq:p171-module-c-powered-arch}.

Finally suppose that \(K\) is trivially valued.  If a normalizing tuple is
not defined over \(K\), pass to the complete isometric nontrivially valued
extension \(F=K((T))\), whose \(T\)-adic absolute value is trivial on
\(K^\times\), and choose the tuple over the infinite field \(F\).  We do
not apply \cite[Proposition~3.8.9]{ChenMoriwaki21AdelicIntersection} across
this change, since the source scope of
\cite[Section~3.8]{ChenMoriwaki21AdelicIntersection} starts with a
nontrivial valuation.  Instead,
\eqref{eq:p171-module-c-deligne-base-change} transports the left side
itself.  Base change of the line-valued tensor, together with
\cite[Remark~1.6.8]{ChenMoriwaki21AdelicIntersection}, identifies
\(F_e\otimes 1\) with the unnormalized resultant polynomial in the
base-changed frame, while
\cite[Remark~3.8.3]{ChenMoriwaki21AdelicIntersection} preserves its epsilon
norm.  The already proved
nontrivially non-Archimedean identity over \(F\) therefore descends as an
unnormalized identity over \(K\).
\cite[Proposition~3.9.3, formula~(3.27)]{ChenMoriwaki21AdelicIntersection}
independently identifies this same epsilon norm with the trivially valued
local intersection number; its proof uses precisely a nontrivial Gauss
extension.  This proves the trivial branch without asserting an unsupported
transport of a chosen normalized local height.

When \(n=0\), all preceding inductions stop at
\cite[Theorem~8.16, formula~(8.11)]{BE21}, and the
line is the positive norm line.  If also \(\dim E_0=1\), the closed
immersion \(Z\hookrightarrow\mathbf P^0_K\) forces
\(Z=\operatorname{Spec}K\); the affine resultant is one linear coordinate,
the projective divisor is empty, and normalization makes its value the
literal unit \(1\).  Thus this endpoint has neither a hidden dualization
nor a missing degree.  This proves
\cref{lem:p171-module-c-field-comparison}.
\end{proof}

\begin{proof}
We prove \cref{thm:p171-module-c-synchronized-comparison}.  The relative
intersection numbers \(d_i\) are constant because they are mixed leading
coefficients of the fibre Hilbert polynomial of a flat projective family
over the connected scheme \(\operatorname{Spec}A\).  Apply
\cref{cor:p171-revised2-veronese-resultants}.  It chooses one Noetherian
integral stage and produces \(\mathcal R_p^{\mathrm{raw}}\); transport by
\(\iota_p\) gives \(\mathcal R_p\), for every sufficiently large \(p\).
On each field fibre it first applies the complete-system theorem
at that stage and only then performs affine monomial substitution.  Thus a
kernel in a monomial evaluation map is harmless and no projective map is
asserted.  The fibre is the positive resultant of \([X_x]_n\).

Ordered multiadditivity gives
\[
 I_{X/A}(M_0^{m_0p},\ldots,M_n^{m_np})
 \simeq I_{X/A}(M_0,\ldots,M_n)^{\otimes\prod_i(m_ip)},
\]
which is \eqref{eq:p171-module-c-line-power}.  The degree of the resultant
in block \(i\) is
\begin{align*}
 (B_0^p\cdots\widehat{B_i^p}\cdots B_n^p)[X_x]_n
 &=\left(\prod_{j\ne i}m_jp\right)d_i\\
 &=\frac{Q_p}{m_ip}d_i,
\end{align*}
which proves \eqref{eq:p171-module-c-powers}, including
\(Q_p=m_0p\) and \(\delta_{0,p}=d_0\) when \(n=0\).

We spell out the analytic localization used next.  The invertible
\(A\)-module underlying \(\mathscr L\) is trivial after inverting one
element \(b\notin\ker(x)\) at any prescribed
\(x\in\mathcal M(A)\).  Since \(\lvert b\rvert_x>0\), there is a rational
Banach analytic neighbourhood of \(x\) on which \(b\) is invertible.  The
chosen algebraic frame and the finitely many coefficients of
\(\mathcal R_p\) then belong to the Banach algebra of that neighbourhood.
These neighbourhoods cover \(\mathcal M(A)\); no finite rational cover and
no global analytic frame is required.

On a local analytic frame, the coefficients in
\eqref{eq:p171-module-c-local-polynomial} are analytic.  If they all
vanished at a point \(x\), the fibre tensor would be the zero polynomial,
contrary to the fibre assertion in
\cref{cor:p171-revised2-veronese-resultants}.  Hence the coefficient tuple
is nowhere zero.  Apply
\cref{thm:hybrid-weighted-philippon-gauss}, with block sizes
\(K_{i,p}\), degrees \(\delta_{i,p}\), and weights
\(w_{p,\alpha}\).  Its definitions are exactly
\eqref{eq:p171-module-c-rho}--\eqref{eq:p171-module-c-phi-arch}.  It proves
choice independence and continuity of \(\Phi_{p,e}\).  Since
\(e'^{\otimes Q_p}=u^{Q_p}e^{\otimes Q_p}\), equality of the two
expressions for \(\mathcal R_p\) gives
\(F_{p,e'}=u^{-Q_p}F_{p,e}\).  Logarithmic homogeneity gives the second
formula in \eqref{eq:p171-module-c-frame-law}, with the displayed minus
sign.

Fix \(x\in U\), put \(K=\mathcal H(x)\), and put
\(\widetilde e_{p,x}=\iota_{p,x}^{-1}(e(x)^{\otimes Q_p})\), a frame of
the raw intersection line
\(I_{X_x/K}(B_{0,x}^p,\ldots,B_{n,x}^p)\).  The definition by transport
gives the literal scalarization
\[
 \mathcal R_{p,x}^{\mathrm{raw}}
 =\widetilde e_{p,x}F_{p,e,x}.
\]
If \(\rho(x)=0\), the dual basis
vector \(U_{i,\beta}\) has norm \(\exp(a_{i,\beta})\).  The epsilon
symmetric quotient norm is exactly
\begin{equation}
 \log\lVert F_{p,e,x}\rVert_\varepsilon
 =\max_\alpha\{\log\lvert c_{p,\alpha}(x)\rvert_x+w_{p,\alpha}\}.
 \label{eq:p171-module-c-exact-gauss}
\end{equation}
For completeness, choose in each tensor word the sorted representative of a
commutative monomial.  This gives a lift with the displayed maximum norm.
Conversely, in any lift the coefficient of a commutative monomial is the
sum of the coefficients of its tensor-word representatives; the
ultrametric inequality bounds its weighted modulus by the norm of the
lift.  Taking the quotient infimum proves equality.  No multinomial
coefficient and no injectivity of the monomial evaluation map is involved.

If \(\rho(x)>0\), the vectors
\(\widehat e_{i,\beta}=\exp(a_{i,\beta}/\rho(x))e_{i,\beta}\)
are orthonormal for the unpowered Hermitian norm.  Their dual coordinates
satisfy
\(U_{i,\beta}=\exp(a_{i,\beta}/\rho(x))\widehat U_{i,\beta}\).
Consequently the coefficient of multiexponent \(\alpha\) in orthonormal
coordinates is multiplied by
\(\exp(w_{p,\alpha}/\rho(x))\), with the positive sign appearing in
\eqref{eq:p171-module-c-arch-polynomial}.  Thus
\cref{lem:p171-module-c-field-comparison}, including its powered real and
complex branches, gives in all four valuation cases
\begin{equation}
 \lambda_x\!\left(\widetilde e_{p,x};
       \widetilde\theta_{0,p,x},\ldots,
       \widetilde\theta_{n,p,x}\right)=\Phi_{p,e}(x).
 \label{eq:p171-module-c-raw-identity}
\end{equation}
The quotient metrics in this formula are continuous semipositive metrics,
so the cited Deligne metric is in the domain of
\cite[Theorem~8.16]{BE21}.

By definition,
\(\widetilde\theta_{i,p,x}=m_ip\,\theta_{i,p,x}\).  Metric
multiadditivity makes \(\iota_{p,x}\) an isometry and gives
\[
 \lambda_x\!\left(\widetilde e_{p,x};
       m_0p\,\theta_{0,p,x},\ldots,m_np\,\theta_{n,p,x}\right)
 =Q_p\lambda_x(e(x);\theta_{0,p,x},\ldots,\theta_{n,p,x}).
\]
Combining this equality with \eqref{eq:p171-module-c-raw-identity} proves
\eqref{eq:p171-module-c-main-identity}.  Notice that the exponent is
\(Q_p\), not a block degree and not its reciprocal.

It remains to prove the uniform approximation, independently of the
coefficient continuity.  Work at a point \(y\in X_x^{\mathrm{an}}\) in a
local frame \(\tau_i\) of \(M_i\), and write
\(s_{ij}=f_{ij}\tau_i^{m_i}\).  At a non-Archimedean point, including a
trivially valued point, the reciprocal quotient norm is
\[
 \max_{\lvert\beta\rvert=p}
 \lvert f_i^\beta(y)\rvert\exp(a_{i,\beta})
 =\left(\max_j\lvert f_{ij}(y)\rvert\exp(a_{ij})\right)^p.
\]
The equality uses the pure powers and multiplicativity.  After taking the
logarithm and dividing by \(m_ip\), it says literally
\(\theta_{i,p,x}=\theta_{i,x}\).

Now let \(\rho=\rho(x)>0\).  Under the canonical unpowered real/complex
analytic identification, choose a complex representative \(y_\infty\) of
\(y\) (the two representatives in the real case are conjugate), and put
\[
 b_j=\log\lvert f_{ij}(y_\infty)\rvert_\infty+a_{ij}/\rho,
 \qquad b_\beta=\sum_j\beta_jb_j.
\]
The Hilbert quotient formula for the declared diagonal monomial basis gives
\begin{equation}
 -\log\lVert\tau_i\rVert_{\theta_{i,p,x}}
 =\frac\rho{2m_ip}\log\sum_{\lvert\beta\rvert=p}\exp(2b_\beta).
 \label{eq:p171-module-c-softmax-formula}
\end{equation}
There are no multinomial coefficients: the norm was declared on the
monomial basis itself, rather than taken as the symmetric-power quotient of
the degree-one Hermitian norm.  Since there are \(K_{i,p}\) terms and the
pure powers give
\(\max_{\lvert\beta\rvert=p}b_\beta=p\max_jb_j\), one has
\[
 \max_\beta b_\beta
 \leq\frac12\log\sum_\beta\exp(2b_\beta)
 \leq\max_\beta b_\beta+\frac12\log K_{i,p}.
\]
The target frame weight is
\[
 -\log\lVert\tau_i\rVert_{\theta_{i,x}}
 =\rho m_i^{-1}\max_jb_j.
\]
Substitution in
\eqref{eq:p171-module-c-softmax-formula} proves
\eqref{eq:p171-module-c-softmax}, including its sign and factor \(1/2\).
The inequality \(\rho\leq1\) follows from the ordinary triangle inequality
\(\lvert2\rvert_x\leq2\).

Change the \(n+1\) normalized Deligne metrics one slot at a time from
\((\theta_i)_i\) to \((\theta_{i,p})_i\).
\cite[Theorem~8.16, formula~(8.9)]{BE21} expresses the change in slot \(i\)
as the integral of \(\theta_{i,p}-\theta_i\) against the positive mixed
measure in the other slots.  At a powered Archimedean fibre this is the
transported formula: \eqref{eq:PN2-r3} contributes one factor \(\rho\),
\eqref{eq:PN1-r3} leaves the unpowered mixed measure, and unpowering
\(\theta_{i,p}-\theta_i\) contributes \(1/\rho\), so no residual power
remains.  Its total mass is \(d_i\), independently of which
already changed semipositive metrics occur there.  Hence
\[
 0\leq\lambda_{p,e}(x)-\lambda_e(x)
 \leq\sum_i d_i\lVert\theta_{i,p,x}-\theta_{i,x}\rVert_\infty
 \leq\sum_i d_i\epsilon_{i,p},
\]
which is \eqref{eq:p171-module-c-deligne-error}.  Equivalently, on the raw
line the exact ledger is
\(\delta_{i,p}(m_ip)=Q_pd_i\), and division by \(Q_p\) gives the same
bound.  This check prevents loss of a Veronese exponent.  When \(n=0\),
the measure in the empty set of other slots has mass
\(d_0=\sum_z\operatorname{length}(A_z)[k_z:K]\), so the same calculation
is literal.  On a real fibre the push-forward convention retains this
mass and adds no factor two.

For fixed \(N_i\),
\(\log K_{i,p}=O(\log p)\), so the right side of
\eqref{eq:p171-module-c-deligne-error} tends to zero independently of
\(x\).  Each \(\lambda_{p,e}=Q_p^{-1}\Phi_{p,e}\) is continuous by the
weighted coefficient theorem.  Therefore \(\lambda_e\) is its uniform
limit and is continuous.  Dividing \eqref{eq:p171-module-c-frame-law} by
\(Q_p\) gives
\(\lambda_{p,ue}=\lambda_{p,e}-\log\lvert u\rvert\); passing to the
uniform limit proves the same law for \(\lambda_e\).  This completes the
proof of \cref{thm:p171-module-c-synchronized-comparison}.
\end{proof}

\begin{proof}
We prove \cref{cor:p171-module-c-finite-level-continuity}.  The source
hypotheses imply finite presentation.  Indeed,
\cite[Section~2]{PS23Hybrid} assumes that \(A\) is integral and \(X\) is
projective of finite type over \(A\).  Flatness and local finite type over an
integral base imply local finite presentation by
\cite[Tag~\texttt{05IN}]{stacks-project}.  Finite type is
quasi-compact, while projective is separated and hence quasi-separated.
The definition in \cite[Tag~\texttt{01TP}]{stacks-project} therefore makes
\(X\to\operatorname{Spec}A\) of finite presentation.

Apply \cref{prop:p171-closed-immersion-finite-level-reduction}.  Its exact
identities \eqref{eq:p171-closed-immersion-dfs-expansion} and
\eqref{eq:p171-closed-immersion-ma-expansion}, followed by the Deligne
change-of-metric formula, express
\eqref{eq:p171-module-c-finite-level-function} as a finite signed sum of
differences of Deligne weights.  The same finite index set and signs work
for every \(x\), and every defining section system in every slot gives a
closed immersion.  The underlying line bundles are therefore ample.

Apply \cref{thm:p171-module-c-synchronized-comparison} separately to each
of those finitely many \((n+1)\)-tuples.  Each Deligne weight in the signed
sum is continuous on \(\mathcal M(A)\); hence so is their finite signed
sum.  The argument also applies for \(n=0\), where it is the difference of
two positive finite-norm weights and includes the Artin lengths and residue
degrees recorded above.  No continuity of the desired integral was used in
constructing the coefficient functions: continuity first came from the
independently verified weighted polynomial theorem, the fieldwise
identification came next, and only then was the uniform limit taken.  Thus
the proof is not circular and establishes only the stated finite-level
conclusion.
\end{proof}
